\documentclass{standalone}
\usepackage{circuitikz}
\usetikzlibrary{calc}
\definecolor{circuitnotemark}{HTML}{FE00FE}
\makeatletter
\pgfdeclareshape{reg3}{
  \anchor{center}{\pgfpointorigin}
  \anchor{text}{\pgfpointorigin}
  \anchor{north}{\pgfpoint{0cm}{0.62cm}}
  \anchor{south}{\pgfpoint{0cm}{-0.62cm}}
  \anchor{east}{\pgfpoint{0.95cm}{0cm}}
  \anchor{west}{\pgfpoint{-0.95cm}{0cm}}
  \anchor{pin 1}{\pgfpoint{-1.35cm}{0cm}}
  \anchor{bpin 1}{\pgfpoint{-0.95cm}{0cm}}
  \anchor{pin 2}{\pgfpoint{0cm}{-1.02cm}}
  \anchor{bpin 2}{\pgfpoint{0cm}{-0.62cm}}
  \anchor{pin 3}{\pgfpoint{1.35cm}{0cm}}
  \anchor{bpin 3}{\pgfpoint{0.95cm}{0cm}}
  \backgroundpath{
    \pgfpathrectanglecorners{\pgfpoint{-0.95cm}{-0.62cm}}{\pgfpoint{0.95cm}{0.62cm}}
    \pgfpathmoveto{\pgfpoint{-1.35cm}{0cm}}\pgfpathlineto{\pgfpoint{-0.95cm}{0cm}}
    \pgfpathmoveto{\pgfpoint{0cm}{-1.02cm}}\pgfpathlineto{\pgfpoint{0cm}{-0.62cm}}
    \pgfpathmoveto{\pgfpoint{1.35cm}{0cm}}\pgfpathlineto{\pgfpoint{0.95cm}{0cm}}
  }
}
\makeatother
\begin{document}
\begin{circuitikz}[american, line width=0.8pt]
\ctikzset{bipoles/length=1.2cm}
\ctikzset{grounds/scale=1.36}
\coordinate (x2y2) at (2,-2);
\coordinate (x2y5) at (2,-8);
\coordinate (x6y3) at (10,-4);
\coordinate (x6y5) at (10,-8);
\coordinate (x8y3) at (14,-4);
\coordinate (x4y2) at (6,-2);
\draw (x2y2) to[esource, l_=$V_{1}$, a^=$5\,\mathrm{V}$] (x2y5); % line 3
\draw (1.919,-4.865) -- ++(0.162,0); % line 3
\draw (2,-4.946) -- ++(0,0.162); % line 3
\draw (1.919,-5.135) -- ++(0.162,0); % line 3
\node[vcc] at (x2y2) {+5V}; % line 4
\node[ground] at (x2y5) {}; % line 5
\node[reg3, draw, font=\scriptsize] (part-U1) at (x6y3) {}; % line 6
\node[font=\scriptsize, anchor=south] at (part-U1.north) {$\mathrm{LM35}$}; % line 6
\node[anchor=south, yshift=9pt] at (part-U1.north) {$U_{1}$}; % line 6
\node[circuitnotemark, font=\tiny, anchor=west, xshift=2pt, inner sep=0] at (part-U1.bpin 1) {X}; % line 6
\node[circuitnotemark, font=\tiny, anchor=west, yshift=4pt, inner sep=0] at (part-U1.bpin 2) {X}; % line 6
\node[circuitnotemark, font=\tiny, anchor=east, xshift=-2pt, inner sep=0] at (part-U1.bpin 3) {X}; % line 6
\draw (x6y5) node[ocirc]{} node[above left]{OUT}; % line 11
\node[ground] at (x8y3) {}; % line 12
\draw (x2y2) -- (x4y2); % line 14
\draw (part-U1.pin 1) -| (x4y2); % line 15
\draw (part-U1.pin 2) -- (x6y5); % line 16
\draw (part-U1.pin 3) -| (x8y3); % line 17
\node[circ] at (x2y2) {};
\node[anchor=south west, inner sep=2pt, yshift=4pt, circuitnotemark, font=\LARGE\bfseries] (circuittitle) at (current bounding box.north west) {X};
\path (circuittitle.south west) rectangle ++(11.552,0.853);
\node[anchor=north east, inner sep=2pt, circuitnotemark, font=\large] (circuitstamp) at ([yshift=-0.593cm]current bounding box.south east) {X};
\path (circuitstamp.north east) rectangle ++(-5.08,-0.593);
\end{circuitikz}
\end{document}
